\documentclass{article}
\usepackage{/globals/de}
\usepackage{/globals/math}
\begin{document}
\section{Logarithmen}
Ein Logarithmus der Zahl $a$ zur Basis $b$ stellt die Frage \textquote{$b$ hoch was ist $a$}.
So gilt
\[
 \log_b{a} = x \Leftrightarrow b^x = a
\]
Die Basis $b$ ist entweder wie oben angegeben, aus dem Kontext offensichtlich, oder es wird eine Kurzform verwendet.
\[
 \ln = \log_\e
 \dtext{und}
 \operatorname{ld} = \log_2
 \dtext{und}
 \operatorname{lg} = \log_{10}
\]

\subsection{Logarithmengesetze}
Manche Ergebnisse sind Konstant.
\[
 \log_a{1} = 0 \dtext{und} \log_a{a} = 1
\]
Eine Multiplikation im Logarithmus kann zu einer Addition von zwei Logarithmen umgewandelt werden.
\[
 \log_a{m*n} = \log_a{m} + \log_a{n}
\]
Es kann immer die Potenz nach vorne gezogen werden.
\[
 \log_a{x^c} = c * \log_a{x}
\]
Mit einer Division durch den Logarithmus einer neuen Basis zur alten Basis kann in eine neue Basis umgewandelt werden.
\[
 \log_{b_2}{x} = \frac{\log_{b_1}{x}}{\log_{b_1}{b_2}}
\]

\end{document}